% =====================================================================
% Phase 1 — Exercices pour aller plus loin (5 exercices bonus)
% Document complémentaire à phase1.tex.
% Compilation : pdflatex -shell-escape phase1-bonus.tex (deux fois pour la TOC)
% Pré-requis  : pygmentize (pip install Pygments) + paquet TeX "minted"
% =====================================================================
\documentclass[11pt,a4paper]{article}

% ---------- Encodage & langue ----------
\usepackage[utf8]{inputenc}
\usepackage[T1]{fontenc}
\usepackage[french]{babel}
\usepackage{lmodern}
\usepackage{microtype}

% ---------- Mise en page ----------
\usepackage[a4paper,margin=2.2cm]{geometry}
\usepackage{parskip}
\usepackage{xcolor}
\usepackage{enumitem}
\usepackage{booktabs}
\usepackage{hyperref}
\hypersetup{
  colorlinks=true,
  linkcolor=blue!60!black,
  urlcolor=blue!60!black,
  citecolor=blue!60!black,
  bookmarksopen=true,
  bookmarksopenlevel=2,
}

% ---------- Boîtes & callouts ----------
\usepackage[most]{tcolorbox}
\tcbset{
  enhanced,
  colback=blue!3,
  colframe=blue!50!black,
  fonttitle=\bfseries,
  arc=2pt, boxrule=0.6pt,
}

% ---------- Code avec minted ----------
\usepackage{minted}
\setminted{
  fontsize=\small,
  linenos=true,
  numbersep=6pt,
  frame=lines,
  framesep=2mm,
  bgcolor=gray!4,
  breaklines=true,
  breakanywhere=true,
  python3=true,
  tabsize=4,
}
\newminted{python}{}

% ---------- TikZ ----------
\usepackage{tikz}
\usetikzlibrary{
  arrows.meta, positioning, shapes.geometric, shapes.misc,
  fit, backgrounds, calc, trees,
  decorations.pathreplacing, decorations.markings,
}

\tikzset{
  node distance=8mm and 12mm,
  block/.style={
    rectangle, rounded corners=2pt, draw=blue!60!black, fill=blue!8,
    minimum width=42mm, minimum height=9mm, align=center, font=\small
  },
  data/.style={
    rectangle, rounded corners=1pt, draw=gray!60, fill=gray!10,
    minimum width=42mm, minimum height=8mm, align=center, font=\footnotesize\ttfamily
  },
  llm/.style={
    rectangle, rounded corners=2pt, draw=red!60!black, fill=red!8,
    minimum width=42mm, minimum height=9mm, align=center, font=\small
  },
  schema/.style={
    rectangle, rounded corners=2pt, draw=teal!60!black, fill=teal!8,
    minimum width=48mm, minimum height=9mm, align=center, font=\small
  },
  parallel/.style={
    rectangle, rounded corners=2pt, draw=violet!60!black, fill=violet!8,
    minimum width=48mm, minimum height=9mm, align=center, font=\small
  },
  predicate/.style={
    diamond, draw=orange!70!black, fill=orange!15, aspect=2.4,
    inner sep=1pt, font=\footnotesize, align=center
  },
  flow/.style={-{Stealth[length=2.2mm]}, thick, draw=blue!60!black},
  flowdata/.style={-{Stealth[length=2.2mm]}, thick, draw=gray!60, dashed},
  label/.style={font=\scriptsize\itshape, fill=white, inner sep=1pt},
  timearrow/.style={-{Stealth[length=2.5mm]}, very thick, draw=black!70},
  bar/.style={rectangle, draw=blue!60!black, fill=blue!20, minimum height=4mm, font=\scriptsize},
}

% ---------- Titre ----------
\title{\textbf{Phase 1 — Exercices pour aller plus loin}\\
\large Cinq exercices bonus avant la Phase 2}
\author{Cours PyTutor2}
\date{}

\begin{document}
\maketitle

\begin{abstract}
\noindent
Ce document complète \texttt{phase1.tex} avec cinq exercices avancés~:
parallélisation à 3 branches, \texttt{.batch()} sur plusieurs inputs,
extension d'un schéma Pydantic, router avec \texttt{RunnableBranch}, et
exécution \texttt{async}. Chaque exercice suit la même structure
qu'en Phase 1 (énoncé en boîte, solution détaillée, vérification dans
LangSmith, récapitulatif).
\end{abstract}

\tableofcontents
\bigskip

% #####################################################################
% #####################################################################
\section{Multi-niveaux 3$\times$ : parallèle vs séquentiel}
% #####################################################################
% #####################################################################

\begin{tcolorbox}[title=Exercice bonus 1. Mesurer le speedup d'un \texttt{RunnableParallel} à 3 branches]
Ajoute un niveau \texttt{"avancé"} à la chaîne multi-niveaux et lance les
\textbf{3 niveaux en parallèle}. Mesure le temps d'exécution
\textbf{séquentiel} (3 \texttt{.invoke()} successifs) et le temps
\textbf{parallèle} (un seul \texttt{.invoke()} sur un \texttt{RunnableParallel}
à 3 clés). Compare et calcule le \emph{speedup}.

\textbf{Contrainte}~: utiliser \texttt{time.perf\_counter()} pour les
mesures (plus précis que \texttt{time.time()}). Afficher le rapport
\texttt{seq\_time / par\_time}.

\bigskip
\noindent\textbf{Objectifs pédagogiques}
\begin{itemize}[leftmargin=*]
  \item Vérifier expérimentalement que \texttt{RunnableParallel} donne une
        latence en $\max(T_i)$ et non $\sum T_i$.
  \item Observer la dégradation du speedup quand le nombre de branches
        augmente (rate limiting côté fournisseur).
  \item Savoir mesurer proprement la latence avec
        \texttt{time.perf\_counter()}.
\end{itemize}
\end{tcolorbox}

\subsection{Solution}

\subsubsection{Architecture}

\begin{center}
\begin{tikzpicture}[node distance=3mm and 6mm]
  % Axe séquentiel
  \draw[timearrow] (0,0) -- (10,0) node[right]{\small temps};
  \node[left] at (-0.2, 0.6) {\scriptsize séquentiel};
  \node[bar, anchor=west, minimum width=18mm] (s1) at (0, 0.6) {beginner};
  \node[bar, anchor=west, minimum width=20mm, fill=orange!30] (s2) at (1.8, 0.6) {intermediate};
  \node[bar, anchor=west, minimum width=22mm, fill=teal!30] (s3) at (3.8, 0.6) {advanced};

  \draw[decorate, decoration={brace, amplitude=2mm}]
       (0, 1.2) -- (6.0, 1.2) node[midway, above, yshift=2mm]{\scriptsize $\sim$6 s};

  % Axe parallèle
  \draw[timearrow] (0,-2.5) -- (10,-2.5);
  \node[left] at (-0.2, -1.4) {\scriptsize parallèle};
  \node[bar, anchor=west, minimum width=18mm] (p1) at (0, -1.4) {beginner};
  \node[bar, anchor=west, minimum width=20mm, fill=orange!30] (p2) at (0, -1.95) {intermediate};
  \node[bar, anchor=west, minimum width=22mm, fill=teal!30] (p3) at (0, -2.5) {advanced};

  \draw[decorate, decoration={brace, amplitude=2mm}]
       (0, -1.0) -- (2.2, -1.0) node[midway, above, yshift=2mm]{\scriptsize $\sim$2 s};
\end{tikzpicture}
\end{center}

Trois branches en parallèle restent à la latence d'\emph{une} seule chaîne.
Le speedup attendu est de $\times 3$ en théorie, $\times 2.5$ à $\times 2.8$
en pratique (overhead de \emph{threading}, rate limits, variations côté
fournisseur).

\subsubsection{Le code complet}

\begin{pythoncode}
import time
from langchain_core.runnables import RunnableParallel
from phase1_explain import explain_chain_for_level


def exercise_1() -> None:
    concept = "les décorateurs Python"

    # --- Construction des 3 chaînes spécialisées ---
    beginner     = explain_chain_for_level("débutant")
    intermediate = explain_chain_for_level("intermédiaire")
    advanced     = explain_chain_for_level("avancé")

    # --- Version SÉQUENTIELLE : 3 invokes successifs ---
    t0 = time.perf_counter()
    seq_results = {
        "beginner":     beginner.invoke({"concept": concept}),
        "intermediate": intermediate.invoke({"concept": concept}),
        "advanced":     advanced.invoke({"concept": concept}),
    }
    seq_time = time.perf_counter() - t0
    print(f"Séquentiel : {seq_time:.2f}s pour 3 niveaux")

    # --- Version PARALLÈLE : un seul invoke sur un RunnableParallel ---
    parallel = RunnableParallel(
        beginner=beginner,
        intermediate=intermediate,
        advanced=advanced,
    )
    t0 = time.perf_counter()
    par_results = parallel.invoke({"concept": concept})
    par_time = time.perf_counter() - t0
    print(f"Parallèle  : {par_time:.2f}s pour 3 niveaux")
    print(f"→ Speedup  : ×{seq_time / par_time:.2f}")
\end{pythoncode}

\subsubsection{Explications pas à pas}

\subsubsection*{Étape 1 — Construire les trois chaînes spécialisées}

\texttt{explain\_chain\_for\_level} est la fabrique qu'on a vue dans le
concept 3 de la Phase 1. Trois appels suffisent à produire trois
\texttt{Runnable}, chacun avec son niveau figé via \texttt{.partial()}.
Pas de duplication de prompt — c'est tout l'intérêt du pattern fabrique.

\subsubsection*{Étape 2 — La mesure séquentielle}

Trois \texttt{.invoke()} successifs s'exécutent en \textbf{série}~: le
thread Python attend la réponse du serveur entre chaque appel. Avec un
modèle moyen, chaque \texttt{.invoke()} prend $\sim$2 s ; total $\sim$6 s.

\subsubsection*{Étape 3 — La mesure parallèle}

\texttt{RunnableParallel(beginner=\dots, intermediate=\dots, advanced=\dots)}
crée un seul \texttt{Runnable} qui, à l'invocation, duplique l'input vers
les 3 branches et les exécute en \textbf{concurrence} (via un pool de
threads). On attend que les trois terminent et on récupère un dict
\texttt{\{"beginner": \dots, "intermediate": \dots, "advanced": \dots\}}.

\subsubsection*{Étape 4 — Le speedup observé}

En pratique, on observe~:
\begin{itemize}[leftmargin=*]
  \item \textbf{Séquentiel} : $\sim$5 à 7 s
  \item \textbf{Parallèle}  : $\sim$2 à 3 s
  \item \textbf{Speedup}    : $\times 2.0$ à $\times 2.8$
\end{itemize}

Le facteur 3 théorique n'est jamais atteint à cause de l'overhead de
threading, des variations de latence côté serveur, et — surtout — des
\emph{rate limits} du fournisseur quand les 3 appels arrivent en même
temps.

\begin{quote}\itshape
\textbf{Coût en tokens} : strictement identique entre séquentiel et
parallèle. La parallélisation ne change rien à la facture, seulement à
la latence perçue.
\end{quote}

\subsection{Vérification dans LangSmith}

L'arbre d'exécution affiche un seul \texttt{RunnableParallel} parent avec
\textbf{3 frères} qui se chevauchent temporellement dans la timeline en
haut du run~:

\begin{center}
\begin{tikzpicture}[every node/.style={font=\footnotesize\ttfamily, anchor=west}]
  \node                              (a) {RunnableParallel<beginner, intermediate, advanced>};
  \node[below=1mm of a, xshift=4mm]  (b) {[beginner] RunnableSequence ← LLM \#1};
  \node[below=1mm of b]              (c) {[intermediate] RunnableSequence ← LLM \#2};
  \node[below=1mm of c]              (d) {[advanced] RunnableSequence ← LLM \#3};
\end{tikzpicture}
\end{center}

Dans la \emph{timeline}, les 3 barres démarrent au même instant
(duplication de l'input) et finissent à $\max(T_i)$. Si l'une d'elles
prend visiblement plus de temps (variations de tokens, retry interne),
le total dépend d'elle.

\subsection*{Récapitulatif}
{\itshape\sloppy
\textit{Ce qu'il faut retenir :}
\begin{enumerate}[leftmargin=*]
  \item \texttt{RunnableParallel(a=\dots, b=\dots, c=\dots)} parallélise
        N chaînes sur le même input. La latence reste à $\max(T_i)$.
  \item \texttt{time.perf\_counter()} est l'outil standard pour mesurer
        une durée en Python (monotone, haute résolution).
  \item Le speedup réel plafonne à $\sim$80\% du facteur théorique à cause
        des rate limits du fournisseur et de l'overhead de threading.
  \item Le coût en tokens reste identique — paralléliser n'est qu'un gain
        de latence, jamais de débit ni d'argent.
\end{enumerate}
\par}

% #####################################################################
% #####################################################################
\section{\texttt{.batch()} vs boucle \texttt{for}}
% #####################################################################
% #####################################################################

\begin{tcolorbox}[title=Exercice bonus 2. Paralléliser plusieurs \emph{inputs} sur la même chaîne]
Utilise \texttt{.batch()} au lieu de \texttt{.invoke()} pour expliquer
une \textbf{liste de concepts} en parallèle. Compare avec une boucle
\texttt{for} qui appelle \texttt{.invoke()} successivement. Calcule le
speedup observé.

\textbf{Contrainte}~: utiliser une chaîne \texttt{intermédiaire} et une
liste d'au moins 4 concepts variés (\texttt{["décorateurs", "métaclasses",
"générateurs", "GIL"]} par exemple).

\bigskip
\noindent\textbf{Objectifs pédagogiques}
\begin{itemize}[leftmargin=*]
  \item Distinguer \texttt{.batch()} (plusieurs inputs, \emph{même} chaîne)
        de \texttt{RunnableParallel} (même input, plusieurs chaînes).
  \item Savoir invoquer une chaîne \texttt{.batch(inputs)} avec une liste
        de dicts.
  \item Contrôler le degré de parallélisme via
        \texttt{config=\{"max\_concurrency": N\}} pour respecter les
        rate limits.
\end{itemize}
\end{tcolorbox}

\subsection{Solution}

\subsubsection{Architecture}

\begin{center}
\begin{tikzpicture}
  \node[data] (inputs) {[\{"concept": "décorateurs"\},\\\ \{"concept": "métaclasses"\},\\\ \{"concept": "générateurs"\},\\\ \{"concept": "GIL"\}]};
  \node[block, below=of inputs] (chain) {chain.batch(\dots)\\\scriptsize (RunnableSequence)};

  \node[llm, below left=10mm and 0mm of chain]    (l1) {invoke \#1};
  \node[llm, below=10mm of chain]                 (l2) {invoke \#2};
  \node[llm, below right=10mm and 0mm of chain]   (l3) {invoke \#3};
  \node[llm, right=of l3]                         (l4) {invoke \#4};

  \node[data, below=22mm of chain] (out) {[Explanation, Explanation, Explanation, Explanation]};

  \draw[flowdata] (inputs) -- (chain);
  \draw[flow]     (chain)  -- (l1) node[label, midway, sloped, above]{en parallèle};
  \draw[flow]     (chain)  -- (l2);
  \draw[flow]     (chain)  -- (l3);
  \draw[flow]     (chain)  -- (l4);
  \draw[flow]     (l1) -- (out);
  \draw[flow]     (l2) -- (out);
  \draw[flow]     (l3) -- (out);
  \draw[flow]     (l4) -- (out);
\end{tikzpicture}
\end{center}

\textbf{Une seule chaîne}, N invocations parallèles. Symétrique mais
distinct de \texttt{RunnableParallel} qui prend N chaînes différentes
sur le même input.

\subsubsection{Le code complet}

\begin{pythoncode}
import time
from phase1_explain import explain_chain_for_level


def exercise_2() -> None:
    concepts = ["les décorateurs", "les métaclasses", "les générateurs", "le GIL"]
    chain = explain_chain_for_level("intermédiaire")

    # --- Version BOUCLE : un invoke après l'autre ---
    t0 = time.perf_counter()
    loop_results = [chain.invoke({"concept": c}) for c in concepts]
    loop_time = time.perf_counter() - t0
    print(f"Boucle for : {loop_time:.2f}s pour {len(concepts)} concepts")

    # --- Version BATCH : un seul appel, parallélisé sous le capot ---
    inputs = [{"concept": c} for c in concepts]
    t0 = time.perf_counter()
    batch_results = chain.batch(inputs)
    batch_time = time.perf_counter() - t0
    print(f"batch()    : {batch_time:.2f}s pour {len(concepts)} concepts")
    print(f"→ Speedup  : ×{loop_time / batch_time:.2f}")

    # Contrôler le degré de parallélisme pour éviter de saturer les rate limits :
    #   chain.batch(inputs, config={"max_concurrency": 5})
\end{pythoncode}

\subsubsection{Explications pas à pas}

\subsubsection*{Étape 1 — Le contrat de \texttt{.batch()}}

\texttt{chain.batch(inputs)} accepte une \textbf{liste} de dicts et
renvoie une liste de résultats \textbf{dans le même ordre} que l'entrée~:

\begin{pythoncode}
inputs = [
    {"concept": "décorateurs"},     # → results[0]
    {"concept": "métaclasses"},     # → results[1]
    {"concept": "générateurs"},     # → results[2]
]
results = chain.batch(inputs)       # même ordre garanti
\end{pythoncode}

Sous le capot, LCEL lance les N invocations en parallèle (threads ou
asyncio selon le contexte) et collecte les résultats.

\subsubsection*{Étape 2 — \texttt{.batch()} vs \texttt{RunnableParallel}}

\begin{center}
\begin{tabular}{p{6.5cm} p{6.5cm}}
\toprule
\textbf{\texttt{chain.batch(inputs)}} & \textbf{\texttt{RunnableParallel(a, b, c)}} \\
\midrule
\textbf{Une seule chaîne}, plusieurs inputs.
& \textbf{N chaînes différentes}, un seul input. \\[1mm]
Entrée : \texttt{list[dict]} (N inputs).
& Entrée : \texttt{dict} unique (dupliqué). \\[1mm]
Sortie : \texttt{list[result]} (même ordre).
& Sortie : \texttt{dict[branch\_name, result]}. \\[1mm]
Usage : explorer N requêtes \emph{indépendantes}.
& Usage : N angles d'analyse d'\emph{une} requête. \\
\bottomrule
\end{tabular}
\end{center}

\begin{quote}\itshape\sloppy
\textbf{Règle mnémotechnique} :
\textbf{plusieurs inputs, même chaîne $\to$ \texttt{.batch()}} ;
\textbf{même input, plusieurs chaînes $\to$ \texttt{RunnableParallel}}.
\end{quote}

\subsubsection*{Étape 3 — Contrôler le parallélisme avec \texttt{max\_concurrency}}

Par défaut, \texttt{.batch(N inputs)} lance les N appels simultanément.
Si N est grand (50, 100\dots), le fournisseur va te \emph{rate limiter}.
Le contrôle se fait via \texttt{config}~:

\begin{pythoncode}
# Limiter à 5 appels concurrents
results = chain.batch(inputs, config={"max_concurrency": 5})
\end{pythoncode}

LCEL exécute alors les inputs en \textbf{vagues de 5}~: 5 partent en
parallèle, on attend qu'ils finissent, puis les 5 suivants. Pour 50
inputs avec \texttt{max\_concurrency=5}, on a 10 vagues séquentielles.

\subsubsection*{Étape 4 — Le speedup observé}

Pour 4 inputs à $\sim$2 s chacun~:
\begin{itemize}[leftmargin=*]
  \item \textbf{Boucle for} : $\sim$8 s
  \item \textbf{batch()}    : $\sim$2 à 3 s
  \item \textbf{Speedup}    : $\times 3$ à $\times 4$
\end{itemize}

\subsection{Vérification dans LangSmith}

Le run racine s'appelle \texttt{RunnableSequence.batch} et contient N
sous-runs frères, chacun étant un \texttt{RunnableSequence} (la chaîne
exécutée sur un input). Dans la timeline, on voit clairement les N barres
parallèles. Le rapport \emph{total\_tokens} sur le run racine agrège les
N appels.

\begin{quote}\itshape
Dans LangSmith, un \texttt{.batch()} produit un seul run parent avec N
enfants frères. Si tu vois N runs racines distincts (au lieu d'un parent
avec N enfants), c'est que tu as utilisé une boucle \texttt{for} et pas
\texttt{.batch()}.
\end{quote}

\subsection*{Récapitulatif}
{\itshape\sloppy
\textit{Ce qu'il faut retenir :}
\begin{enumerate}[leftmargin=*]
  \item \texttt{chain.batch(inputs)} parallélise N invocations d'une
        \textbf{même} chaîne sur une liste de N dicts d'entrée.
  \item Ne pas confondre avec \texttt{RunnableParallel} : batch est
        « plusieurs inputs / même chaîne », RunnableParallel est
        « même input / plusieurs chaînes ».
  \item L'ordre de sortie est garanti identique à l'ordre d'entrée.
  \item \texttt{config=\{"max\_concurrency": N\}} limite le nombre
        d'appels concurrents pour respecter les rate limits.
  \item Speedup typique~: $\times 3$ à $\times 4$ pour 4 inputs, avec
        plafonnement progressif quand N augmente.
\end{enumerate}
\par}

% #####################################################################
% #####################################################################
\section{Schéma Pydantic enrichi avec un champ \texttt{analogy}}
% #####################################################################
% #####################################################################

\begin{tcolorbox}[title=Exercice bonus 3. Étendre un schéma Pydantic sans toucher au prompt]
Ajoute un champ \texttt{analogy: str} au schéma Pydantic de la Phase 1
(analogie avec un objet ou une situation du quotidien). Vérifie que le
modèle remplit ce champ \textbf{sans aucune modification du prompt}.

\textbf{Contrainte}~: la \texttt{description=} du nouveau champ doit
suffire à piloter le contenu. Pas de \texttt{ChatPromptTemplate} modifié,
pas de système re-rédigé. Juste le schéma.

\bigskip
\noindent\textbf{Objectifs pédagogiques}
\begin{itemize}[leftmargin=*]
  \item Mesurer concrètement la \emph{viscosité} (faible) d'un changement
        avec \texttt{with\_structured\_output}.
  \item Comprendre que \texttt{Field(description=\dots)} est un véritable
        levier de \emph{prompt engineering}.
  \item Voir la mise à jour du tool schema dans LangSmith en direct.
\end{itemize}
\end{tcolorbox}

\subsection{Solution}

\subsubsection{Architecture}

\begin{center}
\begin{tikzpicture}[node distance=4mm and 12mm]
  \node[schema] (old) {\texttt{Explanation}\\\scriptsize (Phase 1)};
  \node[data, right=of old] (oldfields)
      {\texttt{concept}\\
       \texttt{intuition}\\
       \texttt{code\_example}\\
       \texttt{pitfall}};

  \node[schema, below=14mm of old, fill=teal!15] (new) {\texttt{ExplanationWithAnalogy}\\\scriptsize (bonus)};
  \node[data, right=of new] (newfields)
      {\texttt{concept}\\
       \texttt{intuition}\\
       \textbf{\texttt{analogy}}\\
       \texttt{code\_example}\\
       \texttt{pitfall}};

  \draw[flow] (old) -- (oldfields);
  \draw[flow] (new) -- (newfields);
  \draw[flowdata, draw=orange!70!black, very thick] (old) -- (new)
       node[label, midway, right]{1 champ ajouté};
\end{tikzpicture}
\end{center}

Une seule ligne change~: l'ajout du champ \texttt{analogy: str = Field(\dots)}
dans la classe. Le reste du code applicatif est strictement identique.

\subsubsection{Le code complet}

\begin{pythoncode}
from pydantic import BaseModel, Field
from phase1_explain import EXPLAIN_PROMPT
from config import get_model


class ExplanationWithAnalogy(BaseModel):
    """Version enrichie : on rajoute une analogie du quotidien."""

    concept: str = Field(description="Le concept expliqué.")
    intuition: str = Field(description="Explication intuitive en 3 phrases.")
    analogy: str = Field(
        description=(
            "Une analogie avec un objet ou une situation du quotidien "
            "(cuisine, sport, bureau...) qui éclaire le concept pour un débutant. "
            "Une seule phrase percutante."
        )
    )
    code_example: str = Field(description="Un exemple de code minimal et commenté.")
    pitfall: str = Field(description="Un piège classique sur ce concept.")


def exercise_3() -> None:
    # Une seule ligne change : le schéma passé à with_structured_output.
    model = get_model().with_structured_output(ExplanationWithAnalogy)
    chain = EXPLAIN_PROMPT.partial(level="débutant") | model

    for concept in ["les décorateurs", "le polymorphisme", "le contexte avec `with`"]:
        result = chain.invoke({"concept": concept})
        print(f"\n--- {concept} ---")
        print(f"Analogie : {result.analogy}")
\end{pythoncode}

\subsubsection{Explications pas à pas}

\subsubsection*{Étape 1 — Étendre la classe Pydantic}

L'ajout d'un champ se fait comme dans n'importe quel \texttt{BaseModel}
Pydantic~: une nouvelle ligne avec un nom, un type, et un
\texttt{Field(description=\dots)}. Le champ devient automatiquement
\textbf{requis} (Pydantic exigera sa présence).

\subsubsection*{Étape 2 — La seule ligne d'application qui change}

\begin{pythoncode}
# Avant
model = get_model().with_structured_output(Explanation)

# Après
model = get_model().with_structured_output(ExplanationWithAnalogy)
\end{pythoncode}

\textbf{C'est tout.} Le prompt système n'a pas bougé. Le \texttt{.partial()}
n'a pas bougé. Le \texttt{|} n'a pas bougé. Le LLM va recevoir le nouveau
tool schema et le remplir intégralement.

\subsubsection*{Étape 3 — La \texttt{description=} fait tout le travail}

Le LLM ne sait pas \emph{a priori} ce qu'on attend d'un champ
\texttt{analogy}. C'est la \texttt{description=} qui le lui dit~:

\begin{quote}\itshape
\textit{« Une analogie avec un objet ou une situation du quotidien
(cuisine, sport, bureau\dots) qui éclaire le concept pour un débutant.
Une seule phrase percutante. »}
\end{quote}

Trois éléments importants dans cette description~:
\begin{itemize}[leftmargin=*]
  \item \textbf{Domaine} : « cuisine, sport, bureau » — donne des exemples
        concrets de registre attendu.
  \item \textbf{Public} : « pour un débutant » — calibre la simplicité.
  \item \textbf{Forme} : « une seule phrase percutante » — limite la
        longueur, force la concision.
\end{itemize}

C'est du \emph{prompt engineering} appliqué directement au schéma.

\subsubsection*{Étape 4 — La force de \texttt{with\_structured\_output}}

C'est cette élasticité qu'on cherche en production~: les contrats de
sortie évoluent (un manager veut un nouveau champ, un produit veut une
nouvelle métadonnée, etc.). Avec un parseur texte classique, il faudrait
rééditer le prompt, gérer les cas d'oubli, écrire du code de fallback.
Avec \texttt{with\_structured\_output}, on ajoute le champ \textbf{et c'est
tout}.

\begin{quote}\itshape
\textbf{Coût} : la description du nouveau champ ajoute quelques tokens
au prompt (le schéma est envoyé à chaque appel). Le coût marginal d'un
champ supplémentaire est de $\sim$30 à 80 tokens — négligeable.
\end{quote}

\subsection{Vérification dans LangSmith}

Dans l'onglet \emph{Tool calls} du \texttt{ChatAnthropic}, on voit le
nouveau tool schema avec 5 champs au lieu de 4. Le champ \texttt{analogy}
apparaît avec sa description complète. L'output du run racine est un
objet \texttt{ExplanationWithAnalogy} dont les 5 champs sont remplis.

\begin{center}
\begin{tikzpicture}[every node/.style={font=\footnotesize\ttfamily, anchor=west}]
  \node                              (a) {ChatAnthropic (structured\_output=ExplanationWithAnalogy)};
  \node[below=1mm of a, xshift=4mm]  (b) {Tools : [\{};
  \node[below=1mm of b, xshift=4mm]  (c) {  name: "ExplanationWithAnalogy",};
  \node[below=1mm of c]              (d) {  parameters: \{ concept, intuition,};
  \node[below=1mm of d, xshift=8mm]  (e) {analogy, code\_example, pitfall \}};
  \node[below=1mm of e, xshift=-12mm](f) {\}]};
\end{tikzpicture}
\end{center}

\begin{quote}\itshape
Modifier la \texttt{description=} d'un champ et relancer : tu vois
immédiatement le changement dans le tool schema envoyé au LLM, et la
sortie générée se met à jour conformément. C'est la boucle la plus rapide
de prompt engineering qu'on puisse avoir.
\end{quote}

\subsection*{Récapitulatif}
{\itshape\sloppy
\textit{Ce qu'il faut retenir :}
\begin{enumerate}[leftmargin=*]
  \item Ajouter un champ à un schéma Pydantic = \textbf{une ligne}, plus
        \textbf{un changement} du nom passé à \texttt{with\_structured\_output}.
  \item La \texttt{description=} d'un \texttt{Field} \textbf{est} une
        partie du prompt. Soigne-la (registre, public, forme).
  \item Le coût marginal d'un champ ajouté est négligeable ($\sim$30 à
        80 tokens supplémentaires).
  \item Dans LangSmith, le tool schema reflète immédiatement les changements
        de description — boucle de prompt engineering ultra-rapide.
\end{enumerate}
\par}

% #####################################################################
% #####################################################################
\section{Router avec \texttt{RunnableBranch} et classifieur LLM}
% #####################################################################
% #####################################################################

\begin{tcolorbox}[title=Exercice bonus 4. Aiguillage dynamique vers 3 sous-chaînes spécialisées]
Crée un \emph{router} qui choisit automatiquement entre 3 sous-chaînes
selon que le concept Python est de catégorie \textbf{syntaxe} (def, lambda,
comprehensions\dots), \textbf{paradigme} (POO, fonctionnel,
polymorphisme\dots) ou \textbf{écosystème} (pip, venv, pytest\dots).
Utilise une mini-chaîne LLM de classification en amont pour décider du
routage.

\textbf{Contrainte}~: tout doit rester un seul \texttt{Runnable} LCEL
pour préserver streaming, batching, et tracing LangSmith unifié.

\bigskip
\noindent\textbf{Objectifs pédagogiques}
\begin{itemize}[leftmargin=*]
  \item Maîtriser le \textbf{pattern router}~: classifieur LLM + sous-chaînes
        spécialisées + \texttt{RunnableBranch}.
  \item Utiliser \texttt{Literal[\dots]} dans un schéma Pydantic pour
        garantir un ensemble fermé de catégories.
  \item Combiner \texttt{RunnablePassthrough.assign} et
        \texttt{RunnableBranch} pour enrichir puis brancher.
  \item Repérer dans LangSmith les \textbf{trois appels LLM} qu'un router
        engendre par requête (et le surcoût associé).
\end{itemize}
\end{tcolorbox}

\subsection{Solution}

\subsubsection{Architecture}

\begin{center}
\begin{tikzpicture}
  \node[data] (in) {\{"concept": "lambda"\}};
  \node[block, below=of in] (assign) {RunnablePassthrough.assign\\(category=classifier)};
  \node[data, below=of assign] (enriched) {\{"concept": "lambda",\\\ "category": "syntaxe"\}};
  \node[block, below=of enriched] (branch) {RunnableBranch\\(3 cas + défaut)};
  \node[llm, below=of branch] (out) {réponse spécialisée\\(AIMessage)};

  \draw[flowdata] (in)       -- (assign);
  \draw[flow]     (assign)   -- (enriched) node[label, midway, right]{enrichit le dict};
  \draw[flowdata] (enriched) -- (branch);
  \draw[flow]     (branch)   -- (out)      node[label, midway, right]{choisit la sous-chaîne};
\end{tikzpicture}
\end{center}

\subsubsection*{Zoom sur le \texttt{RunnableBranch}}

\begin{center}
\begin{tikzpicture}[node distance=5mm and 12mm]
  \node[data] (in) {\{"concept": "lambda",\\\ "category": "syntaxe"\}};
  \node[predicate, right=18mm of in] (p1) {category\\=="syntaxe"?};
  \node[predicate, below=10mm of p1] (p2) {category\\=="paradigme"?};
  \node[block, below=10mm of p2] (def) {ecosysteme\_chain\\\scriptsize (défaut)};
  \node[llm, right=of p1] (c1) {syntaxe\_chain};
  \node[llm, right=of p2] (c2) {paradigme\_chain};

  \draw[flowdata] (in) -- (p1);
  \draw[flow]     (p1) -- node[label, above]{vrai} (c1);
  \draw[flow]     (p1) -- node[label, right]{faux} (p2);
  \draw[flow]     (p2) -- node[label, above]{vrai} (c2);
  \draw[flow]     (p2) -- node[label, right]{faux} (def);
\end{tikzpicture}
\end{center}

\subsubsection{Le code complet}

\begin{pythoncode}
from typing import Literal
from pydantic import BaseModel, Field
from langchain_core.prompts import ChatPromptTemplate
from langchain_core.runnables import (
    RunnableBranch, RunnableLambda, RunnablePassthrough,
)
from config import get_model


class CategoryResult(BaseModel):
    """Catégorie d'un concept Python pour router vers la bonne sous-chaîne."""
    category: Literal["syntaxe", "paradigme", "ecosysteme"] = Field(
        description=(
            "Catégorise le concept :\n"
            " - 'syntaxe'    : éléments du langage (def, lambda, comprehensions, with, yield...)\n"
            " - 'paradigme'  : concepts de design (POO, fonctionnel, héritage, polymorphisme...)\n"
            " - 'ecosysteme' : outils/bibliothèques (pip, venv, pytest, numpy, requests...)"
        )
    )


def build_classifier():
    """Mini-chaîne qui catégorise un concept en 3 classes."""
    classify_prompt = ChatPromptTemplate.from_messages([
        ("system", "Tu classes des concepts Python en 3 catégories. "
                   "Ne réponds que par la catégorie."),
        ("human",  "Concept à catégoriser : {concept}"),
    ])
    return (
        classify_prompt
        | get_model().with_structured_output(CategoryResult)
        | RunnableLambda(lambda r: r.category)
    )


def build_specialized_chain(focus: str):
    """Construit une sous-chaîne avec un prompt système orienté `focus`."""
    prompt = ChatPromptTemplate.from_messages([
        ("system",
         "Tu es un assistant Python. Le concept à expliquer relève de la "
         f"catégorie '{focus}'. Adapte ton explication en conséquence :\n" +
         {
             "syntaxe":    "insiste sur la syntaxe exacte, montre un mini exemple et un anti-exemple.",
             "paradigme":  "insiste sur la philosophie, le 'pourquoi', les implications de design.",
             "ecosysteme": "insiste sur l'usage pratique, l'installation, les commandes shell typiques.",
         }[focus]),
        ("human", "Explique : {concept}"),
    ])
    return prompt | get_model()


def build_router():
    """Le routeur complet : classifie puis branche."""
    classifier        = build_classifier()
    syntaxe_chain     = build_specialized_chain("syntaxe")
    paradigme_chain   = build_specialized_chain("paradigme")
    ecosysteme_chain  = build_specialized_chain("ecosysteme")

    branch = RunnableBranch(
        (lambda x: x["category"] == "syntaxe",   syntaxe_chain),
        (lambda x: x["category"] == "paradigme", paradigme_chain),
        ecosysteme_chain,  # défaut : si rien ne matche
    )

    return RunnablePassthrough.assign(category=classifier) | branch
\end{pythoncode}

\subsubsection{Explications pas à pas}

\subsubsection*{Étape 1 — Le vocabulaire de routage (\texttt{CategoryResult})}

Avant de router, il faut un \textbf{ensemble fermé} de catégories. Le
\texttt{Literal["syntaxe", "paradigme", "ecosysteme"]} est crucial~:
il \emph{garantit} que le LLM ne pourra pas renvoyer \texttt{"truc"} ou
\texttt{"je sais pas"}. La \texttt{description=} du champ sert de
mini-prompt qui explique au modèle \emph{comment} classer.

\subsubsection*{Étape 2 — Le classifieur}

Trois maillons en pipeline~:
\begin{enumerate}[leftmargin=*]
  \item \texttt{classify\_prompt}~: \texttt{\{"concept": "lambda"\}}
        $\rightarrow$ messages formatés.
  \item \texttt{model.with\_structured\_output(CategoryResult)}~:
        messages $\rightarrow$ objet \texttt{CategoryResult(category="syntaxe")}.
  \item \texttt{RunnableLambda(lambda r: r.category)}~: objet $\rightarrow$
        chaîne \texttt{"syntaxe"}.
\end{enumerate}

Le maillon 3 n'est pas cosmétique~: il rend la sortie utilisable comme
\textbf{valeur de comparaison} dans les prédicats du
\texttt{RunnableBranch}.

\subsubsection*{Étape 3 — Les sous-chaînes spécialisées}

\texttt{build\_specialized\_chain(focus)} est une \emph{fabrique}~: on
l'appelle 3 fois avec des \texttt{focus} différents et on obtient 3
\texttt{Runnable} distincts, chacun avec sa personnalité dans le system
prompt.

\begin{quote}\itshape
La spécialisation se fait à la \textbf{construction}, pas à l'appel.
\texttt{build\_specialized\_chain("syntaxe")} renvoie un \texttt{Runnable}
prêt à l'emploi qui ne sait répondre que dans le style « syntaxe ».
\end{quote}

\subsubsection*{Étape 4 — Le \texttt{RunnableBranch}}

\texttt{RunnableBranch} est le \texttt{if/elif/else} de LCEL. Trois
points clés~:

\begin{itemize}[leftmargin=*]
  \item Les prédicats sont évalués \textbf{dans l'ordre}, le premier qui
        renvoie \texttt{True} gagne (court-circuit).
  \item Si \emph{aucun} prédicat ne matche, la chaîne par défaut s'exécute
        — elle est \textbf{obligatoire}.
  \item La sous-chaîne sélectionnée reçoit \textbf{l'input complet} (le
        dict entier), pas juste le résultat du prédicat.
\end{itemize}

\subsubsection*{Étape 5 — L'enrichissement avant branchement}

Le \texttt{RunnableBranch} a besoin que \texttt{category} \textbf{soit
déjà dans l'input} pour que ses prédicats puissent le lire. C'est le rôle
de \texttt{RunnablePassthrough.assign}~:

\begin{quote}\itshape
\texttt{RunnablePassthrough.assign(category=classifier)} appelle le
classifieur et ajoute le résultat sous la clé \texttt{category}, sans
toucher aux clés existantes. C'est l'équivalent d'un
\texttt{SELECT *, expr AS col} en SQL.
\end{quote}

\subsection{Vérification dans LangSmith}

\subsubsection*{L'arbre d'exécution}

Pour le concept \texttt{"lambda"}, on observe \textbf{3 appels LLM}~:

\begin{center}
\begin{tikzpicture}[every node/.style={font=\footnotesize\ttfamily, anchor=west}]
  \node                              (a) {RunnableSequence (router)};
  \node[below=1mm of a, xshift=4mm]  (b) {RunnableAssign<category>};
  \node[below=1mm of b, xshift=4mm]  (c) {classifier ← LLM \#1 (structured\_output=CategoryResult)};
  \node[below=1mm of c]              (d) {  → "syntaxe"};
  \node[below=1mm of d, xshift=-4mm] (e) {RunnableBranch};
  \node[below=1mm of e, xshift=4mm]  (f) {branch[0]: syntaxe\_chain};
  \node[below=1mm of f, xshift=4mm]  (g) {ChatAnthropic ← LLM \#2 (réponse spécialisée)};
\end{tikzpicture}
\end{center}

\subsubsection*{Les quatre choses à regarder}

\begin{enumerate}[leftmargin=*]
  \item \textbf{L'arbre lui-même}~: tu vois \emph{visuellement} le passage
        par le classifieur puis la branche choisie. C'est l'aiguillage
        dynamique en image.
  \item \textbf{Premier \texttt{ChatAnthropic}} (classifieur)~: onglet
        \emph{Output} → l'objet \texttt{CategoryResult} et le tool-call
        que Claude a utilisé.
  \item \textbf{Deuxième \texttt{ChatAnthropic}} (réponse spécialisée)~:
        le system prompt commence par « Tu es un assistant Python\dots
        catégorie 'syntaxe' ». Précieux pour debugger.
  \item \textbf{Onglet \emph{Costs}} : tokens et coût agrégés sur les 2
        (ou 3, si on garde le \texttt{debug\_chain}) appels. Permet de
        chiffrer le surcoût du pattern router.
\end{enumerate}

\begin{quote}\itshape
Compare deux runs dans l'UI : \texttt{"lambda"} déclenche
\texttt{syntaxe\_chain}, \texttt{"pytest"} déclenche
\texttt{ecosysteme\_chain}. On voit littéralement \textbf{la branche
choisie} changer selon le concept.
\end{quote}

\subsection*{Récapitulatif}
{\itshape\sloppy
\textit{Ce qu'il faut retenir :}
\begin{enumerate}[leftmargin=*]
  \item \textbf{Pattern router} = classifieur LLM + sous-chaînes
        spécialisées + \texttt{RunnableBranch}.
  \item \texttt{with\_structured\_output} + \texttt{Literal} = garantie
        que la sortie de classification est dans un ensemble fermé, donc
        routable sans validation supplémentaire.
  \item \texttt{RunnablePassthrough.assign(\dots)} enrichit le dict de
        contexte au fil de la pipeline, sans perdre les clés précédentes.
  \item Le tout reste un \textbf{seul \texttt{Runnable}}~:
        \texttt{router.invoke}, \texttt{router.batch},
        \texttt{router.astream} fonctionnent — et LangSmith trace l'arbre
        complet en un seul run.
  \item Surcoût~: \textbf{2 appels LLM par requête} (classification +
        réponse), au lieu d'un seul avec un prompt méga-universel.
        Acceptable pour gagner en qualité, mesurable dans LangSmith.
\end{enumerate}
\par}

% #####################################################################
% #####################################################################
\section{Async avec \texttt{asyncio.gather} vs \texttt{RunnableParallel}}
% #####################################################################
% #####################################################################

\begin{tcolorbox}[title=Exercice bonus 5. La voie async pure et sa comparaison avec RunnableParallel]
Réécris la démo multi-niveaux en \textbf{async} avec \texttt{await
chain.ainvoke(\dots)} et \texttt{asyncio.gather(\dots)}. Compare les
durées avec la version \texttt{RunnableParallel} (synchrone et asynchrone).

\textbf{Contrainte}~: utiliser \texttt{asyncio.run()} comme point d'entrée
pour la fonction \texttt{async}, et \texttt{await} sur \texttt{.ainvoke()}
(la variante asynchrone de \texttt{.invoke()}).

\bigskip
\noindent\textbf{Objectifs pédagogiques}
\begin{itemize}[leftmargin=*]
  \item Connaître les méthodes \texttt{.ainvoke()}, \texttt{.abatch()},
        \texttt{.astream()} — les variantes async de tous les
        \texttt{Runnable}.
  \item Comprendre que \texttt{RunnableParallel} utilise asyncio
        \textbf{sous le capot}~: la version async et la version sync
        donnent des temps quasi identiques.
  \item Choisir entre \texttt{asyncio.gather} (contrôle explicite) et
        \texttt{RunnableParallel} (composition implicite).
\end{itemize}
\end{tcolorbox}

\subsection{Solution}

\subsubsection{Architecture}

\begin{center}
\begin{tikzpicture}[node distance=4mm and 14mm]
  \node[data] (in) {\{"concept": "générateurs"\}};

  \node[block, below left=14mm and 14mm of in] (gather)
       {asyncio.gather\\\scriptsize (3 ainvoke en parallèle)};
  \node[parallel, below right=14mm and 14mm of in] (par)
       {RunnableParallel.ainvoke\\\scriptsize (3 chaînes en parallèle)};

  \node[llm, below=of gather] (g_out) {résultats : tuple\\(beg\_r, int\_r, adv\_r)};
  \node[llm, below=of par]    (p_out) {résultats : dict\\\{"beginner": \dots, \dots\}};

  \draw[flowdata] (in) -- (gather);
  \draw[flowdata] (in) -- (par);
  \draw[flow] (gather) -- (g_out);
  \draw[flow] (par) -- (p_out);

  \node[font=\scriptsize\itshape, below=24mm of in] {même latence, même coût, deux styles d'écriture};
\end{tikzpicture}
\end{center}

Les deux approches s'appuient sur \textbf{le même runtime asyncio}~:
event loop unique, 3 coroutines en concurrence, attente non bloquante.
La différence est purement stylistique.

\subsubsection{Le code complet}

\begin{pythoncode}
import asyncio
import time
from langchain_core.runnables import RunnableParallel
from phase1_explain import explain_chain_for_level


async def exercise_5() -> None:
    concept      = "les générateurs Python"
    beginner     = explain_chain_for_level("débutant")
    intermediate = explain_chain_for_level("intermédiaire")
    advanced     = explain_chain_for_level("avancé")

    # --- Version asyncio.gather pure ---
    t0 = time.perf_counter()
    beg_r, int_r, adv_r = await asyncio.gather(
        beginner.ainvoke({"concept": concept}),
        intermediate.ainvoke({"concept": concept}),
        advanced.ainvoke({"concept": concept}),
    )
    gather_time = time.perf_counter() - t0
    print(f"asyncio.gather    : {gather_time:.2f}s")

    # --- Version RunnableParallel avec .ainvoke (async aussi) ---
    parallel = RunnableParallel(
        beginner=beginner, intermediate=intermediate, advanced=advanced,
    )
    t0 = time.perf_counter()
    results = await parallel.ainvoke({"concept": concept})
    parallel_time = time.perf_counter() - t0
    print(f"RunnableParallel  : {parallel_time:.2f}s")

    # --- Version synchrone pour référence ---
    t0 = time.perf_counter()
    parallel.invoke({"concept": concept})
    sync_time = time.perf_counter() - t0
    print(f"Sync (RunnablePar): {sync_time:.2f}s")


# Point d'entrée : asyncio.run lance la coroutine
if __name__ == "__main__":
    asyncio.run(exercise_5())
\end{pythoncode}

\subsubsection{Explications pas à pas}

\subsubsection*{Étape 1 — Le préfixe \texttt{a} : la version async de tout}

LCEL définit une variante async pour chaque méthode synchrone d'un
\texttt{Runnable}~:

\begin{center}
\begin{tabular}{ll}
\toprule
\textbf{Synchrone} & \textbf{Asynchrone} \\
\midrule
\texttt{.invoke(input)}      & \texttt{.ainvoke(input)} \\
\texttt{.batch(inputs)}      & \texttt{.abatch(inputs)} \\
\texttt{.stream(input)}      & \texttt{.astream(input)} \\
\bottomrule
\end{tabular}
\end{center}

Les versions async retournent des \texttt{coroutine}, à attendre avec
\texttt{await} dans un contexte \texttt{async def}.

\subsubsection*{Étape 2 — \texttt{asyncio.gather} pur}

\texttt{asyncio.gather(*coros)} prend N coroutines et les exécute en
\textbf{concurrence} sur l'event loop. \texttt{await} sur le résultat
attend que toutes terminent et renvoie un \textbf{tuple} des résultats
dans l'ordre des arguments~:

\begin{pythoncode}
beg_r, int_r, adv_r = await asyncio.gather(
    beginner.ainvoke({"concept": concept}),     # → beg_r
    intermediate.ainvoke({"concept": concept}), # → int_r
    advanced.ainvoke({"concept": concept}),     # → adv_r
)
\end{pythoncode}

C'est l'équivalent \emph{asynchrone} d'un \texttt{ThreadPoolExecutor}~:
le main thread ne fait rien pendant que les 3 appels HTTP sont en vol.

\subsubsection*{Étape 3 — \texttt{RunnableParallel} utilise asyncio sous le capot}

Quand on appelle \texttt{parallel.invoke(\dots)} (synchrone), LCEL crée
en interne un event loop, lance les branches comme des coroutines, et
attend le tout. Quand on appelle \texttt{parallel.ainvoke(\dots)}, on
réutilise l'event loop existant — plus efficace, surtout dans une app
web déjà en async.

\begin{quote}\itshape
\textbf{Conséquence pratique} : \texttt{parallel.invoke()} et
\texttt{await parallel.ainvoke()} donnent des temps \emph{identiques} à
la milliseconde près. Le choix est dicté par le contexte d'appel
(synchrone ou asynchrone), pas par la performance.
\end{quote}

\subsubsection*{Étape 4 — Quand utiliser \texttt{gather} vs \texttt{RunnableParallel}}

\begin{center}
\begin{tabular}{p{6.5cm} p{6.5cm}}
\toprule
\textbf{\texttt{asyncio.gather}} & \textbf{\texttt{RunnableParallel}} \\
\midrule
Bibliothèque standard Python.   & Primitive LCEL. \\[1mm]
Contrôle explicite du flot de coroutines.
& Composition déclarative dans une chaîne. \\[1mm]
Bien pour du code applicatif~: routes FastAPI, boucle de chat, jobs.
& Bien pour la \emph{composition} dans un \texttt{Runnable} (qui sera lui-même streamé, batché, tracé). \\[1mm]
Pas de tracing LangSmith par défaut (faut envelopper manuellement).
& Tracing LangSmith automatique. \\
\bottomrule
\end{tabular}
\end{center}

\begin{quote}\itshape\sloppy
\textbf{Règle pratique} : \texttt{RunnableParallel} dans les
\textbf{chaînes} (parce que c'est composable et tracé) ;
\texttt{asyncio.gather} dans le \textbf{code applicatif} (parce que
c'est explicite et standard Python).
\end{quote}

\subsubsection*{Étape 5 — Le point d'entrée \texttt{asyncio.run}}

Une fonction \texttt{async def} ne s'exécute pas avec un appel direct~:
elle renvoie une coroutine non encore lancée. \texttt{asyncio.run(coro)}
crée un event loop, exécute la coroutine, et ferme proprement.

\begin{pythoncode}
async def exercise_5() -> None:
    ...                          # contient des await

# Dans le main, on emballe avec asyncio.run
if __name__ == "__main__":
    asyncio.run(exercise_5())
\end{pythoncode}

\subsection{Vérification dans LangSmith}

Les deux versions (\texttt{gather} et \texttt{RunnableParallel.ainvoke})
produisent des arbres très similaires dans LangSmith. La différence~:

\begin{itemize}[leftmargin=*]
  \item \texttt{RunnableParallel.ainvoke}~: un \emph{seul} run racine
        \texttt{RunnableParallel} avec 3 enfants (un par branche).
  \item \texttt{asyncio.gather} avec 3 \texttt{.ainvoke()} séparés~:
        \textbf{3 runs racines distincts}, chacun étant la trace d'un
        \texttt{RunnableSequence}. Pas d'agrégation automatique.
\end{itemize}

Pour avoir un seul run agrégé en \texttt{gather}, il faut envelopper
manuellement avec \texttt{trace} de \texttt{langsmith}~:

\begin{pythoncode}
from langsmith import trace

with trace("multi-level parallel"):
    beg_r, int_r, adv_r = await asyncio.gather(...)
\end{pythoncode}

\begin{quote}\itshape
La traçabilité \textbf{automatique} est un argument fort pour
\texttt{RunnableParallel}. Dès qu'on sort de LCEL, le tracing devient
manuel.
\end{quote}

\subsection*{Récapitulatif}
{\itshape\sloppy
\textit{Ce qu'il faut retenir :}
\begin{enumerate}[leftmargin=*]
  \item Tout \texttt{Runnable} expose une variante async~:
        \texttt{.ainvoke}, \texttt{.abatch}, \texttt{.astream}. Le préfixe
        \texttt{a} = \emph{async}.
  \item \texttt{asyncio.gather(*coros)} lance N coroutines en concurrence
        et renvoie un tuple des résultats dans l'ordre.
  \item \texttt{RunnableParallel} utilise asyncio sous le capot —
        \texttt{.invoke()} et \texttt{await .ainvoke()} ont les mêmes
        performances.
  \item Règle pratique~: \texttt{RunnableParallel} dans les chaînes (pour
        la composition et le tracing), \texttt{asyncio.gather} dans le
        code applicatif (pour le contrôle explicite).
  \item \texttt{asyncio.run(coro)} est le point d'entrée pour appeler une
        coroutine depuis un contexte synchrone (script, \texttt{if
        \_\_name\_\_ == "\_\_main\_\_"}).
\end{enumerate}
\par}

\end{document}
